%% kapa: exam=1 round=32 year=2021 period=2 form=A subject=accounting subject_ko=회계학 seq=11 number=51
%% source: Q-Net 감정평가사 기출문제 (cst00309) artlSeq=5208397 file=extracted/1008/5208397/A형/2021년도 제32회 감정평가사 1차 2교시 A형.pdf
%% conversion: kordoc 4.24.51 -> tools/gen1.py (auto); review=proofread-sonnet; flags=none
\begin{question}{51}
고객과의 계약에서 생기는 수익에 관한 설명으로 옳지 않은 것은?
\begin{choices}
  \item 거래가격을 산정하기 위해서는 계약 조건과 기업의 사업 관행을 참고하며, 거래가격에는 제삼자를 대신해서 회수한 금액은 제외한다.
  \item 고객과의 계약에서 약속한 대가는 고정금액, 변동금액 또는 둘 다를 포함할 수 있다.
  \item 변동대가의 추정이 가능한 경우, 계약에서 가능한 결과치가 두 가지뿐일 경우에는 기댓값이 변동대가의 적절한 추정치가 될 수 있다.
  \item 기업이 받을 권리를 갖게 될 변동대가(금액)에 미치는 불확실성의 영향을 추정할 때에는 그 계약 전체에 하나의 방법을 일관되게 적용한다.
  \item 고객에게서 받은 대가의 일부나 전부를 고객에게 환불할 것으로 예상하는 경우에는 환불부채를 인식한다.
\end{choices}
\end{question}
\answer{3}
